% ============================================================
% DEUTSCH LERNEN IN MANNHEIM — Premium Lernbuch
% LaTeX Template — KDP-ready (6x9 Zoll, 300 DPI)
% ============================================================
\documentclass[12pt,twoside,openright]{memoir}

% --- PAKETE ---
\usepackage{fontspec}           % Professionelle Fonts (XeLaTeX)
\usepackage{microtype}          % Optimale Typografie, Silbentrennung
\usepackage{geometry}           % Seitenlayout
\usepackage{graphicx}           % KI-Bilder importieren
\usepackage{xcolor}             % Farben (Gold, Schwarz)
\usepackage{titlesec}           % Kapitelüberschriften anpassen
\usepackage{hyperref}           % Interaktives Inhaltsverzeichnis
\usepackage{polyglossia}        % Mehrsprachig (DE + UA)
\usepackage{multicol}           % Zweispaltig für Glossar
\usepackage{booktabs}           % Schöne Tabellen
\usepackage{enumitem}           % Listen anpassen
\usepackage{tcolorbox}          % Bunte Lernkästen
\usepackage{fancyhdr}           % Kopf-/Fußzeilen

% --- SEITENLAYOUT (Amazon KDP Standard: 6x9 Zoll) ---
\geometry{
  paperwidth=6in,
  paperheight=9in,
  inner=2.5cm,
  outer=2cm,
  top=2.5cm,
  bottom=2.5cm,
  includeheadfoot
}

% --- FARBEN ---
\definecolor{gold}{RGB}{212, 175, 55}
\definecolor{darkbg}{RGB}{26, 26, 26}
\definecolor{lightgold}{RGB}{201, 162, 39}
\definecolor{softblue}{RGB}{70, 130, 180}
\definecolor{warmgray}{RGB}{245, 243, 238}

% --- FONTS ---
\setmainfont{Georgia}
\setsansfont{Arial}
\setmonofont{Courier New}

% --- SPRACHEN ---
\setmainlanguage{german}
\setotherlanguage{ukrainian}

% --- HYPERREF KONFIGURATION ---
\hypersetup{
  colorlinks=true,
  linkcolor=gold,
  urlcolor=softblue,
  pdftitle={Deutsch lernen in Mannheim},
  pdfauthor={German Cards Ukraine},
  pdfsubject={Deutsch Lernbuch für Anfänger},
  pdfkeywords={Deutsch, Mannheim, Ukraine, Anfänger, Lernbuch}
}

% --- KAPITELÜBERSCHRIFTEN ---
\titleformat{\chapter}[display]
  {\normalfont\sffamily\Huge\bfseries\color{gold}}
  {\chaptertitlename\ \thechapter}
  {20pt}
  {\Huge}
\titlespacing*{\chapter}{0pt}{50pt}{40pt}

\titleformat{\section}
  {\normalfont\sffamily\Large\bfseries\color{gold}}
  {\thesection}{1em}{}

\titleformat{\subsection}
  {\normalfont\sffamily\large\bfseries}
  {\thesubsection}{1em}{}

% --- LERNKÄSTEN ---
\tcbuselibrary{skins,breakable}

\newtcolorbox{lernkasten}[1][]{
  colback=warmgray,
  colframe=gold,
  fonttitle=\sffamily\bfseries,
  title=#1,
  boxrule=1.5pt,
  arc=4pt
}

\newtcolorbox{uebungskasten}[1][]{
  colback=white,
  colframe=softblue,
  fonttitle=\sffamily\bfseries,
  title=#1,
  boxrule=1.5pt,
  arc=4pt
}

\newtcolorbox{glossarkasten}[1][]{
  colback=warmgray,
  colframe=lightgold,
  fonttitle=\sffamily\bfseries\large,
  title=#1,
  boxrule=1pt
}

% --- KOPF- UND FUSSZEILEN ---
\pagestyle{fancy}
\fancyhf{}
\fancyhead[LE]{\small\sffamily\color{lightgold}Deutsch lernen in Mannheim}
\fancyhead[RO]{\small\sffamily\color{lightgold}\leftmark}
\fancyfoot[C]{\small\sffamily\thepage}
\renewcommand{\headrulewidth}{0.4pt}
\renewcommand{\headrule}{\hbox to\headwidth{\color{gold}\leaders\hrule height \headrulewidth\hfill}}

% --- ZEILENABSTAND ---
\setlength{\parskip}{0.5em}
\setlength{\parindent}{1.5em}

% ============================================================
% DOKUMENT BEGINNT
% ============================================================
\begin{document}

% --- TITELSEITE ---
\begin{titlingpage}
  \begin{center}
    \vspace*{2cm}

    % Logo / KI-Bild Cover
    % \includegraphics[width=0.8\textwidth]{images/cover.png}

    \vspace{1cm}
    {\color{gold}\rule{\linewidth}{3pt}}
    \vspace{0.5cm}

    {\fontsize{36}{42}\selectfont\sffamily\bfseries\color{darkbg}
    DEUTSCH LERNEN\\IN MANNHEIM}

    \vspace{0.5cm}
    {\color{gold}\rule{\linewidth}{3pt}}
    \vspace{1cm}

    {\Large\sffamily Premium Lernbuch für Anfänger}

    \vspace{0.3cm}
    {\large\sffamily\color{lightgold}З українськими перекладами}

    \vfill

    {\sffamily\large German Cards Ukraine}\\
    {\sffamily\small\color{lightgold}Telegram: @germancardsua}

    \vspace{1cm}
  \end{center}
\end{titlingpage}

% --- INHALTSVERZEICHNIS ---
\tableofcontents
\newpage

% ============================================================
% KAPITEL 1: EINLEITUNG
% ============================================================
\chapter{Willkommen in Mannheim!}

% Kapitel-Illustration (KI-Bild)
\begin{figure}[h]
  \centering
  % \includegraphics[width=0.7\textwidth]{images/mannheim-wasserturm.png}
  \fbox{\rule{0pt}{4cm}\rule{8cm}{0pt}} % Platzhalter
  \caption{Der Wasserturm — das Wahrzeichen Mannheims}
\end{figure}

\begin{lernkasten}[Herzlich willkommen! Ласкаво просимо!]
Dieses Buch wurde speziell für Menschen aus der Ukraine entwickelt, die in Mannheim leben und Deutsch lernen möchten. Jede Lektion enthält:
\begin{itemize}[nosep]
  \item Wichtige Wörter auf Deutsch
  \item Ukrainische Übersetzungen
  \item Typische Situationen aus dem Alltag
  \item Übungen zum Üben
\end{itemize}
\end{lernkasten}

\section{Wie nutze ich dieses Buch?}

Jedes Kapitel ist aufgebaut wie ein echter Schritt in deinem neuen Leben in Mannheim. Du lernst die Sprache so, wie sie wirklich gebraucht wird — nicht aus der Theorie, sondern aus dem echten Alltag.

\vspace{1em}
\begin{uebungskasten}[Dein erstes Deutsch]
  Sag laut nach: \textbf{``Guten Morgen!''} \\
  Auf Ukrainisch: \textit{Доброго ранку!}
\end{uebungskasten}

% ============================================================
% KAPITEL 2: GRUNDLAGEN
% ============================================================
\chapter{Grundlagen — Begrüßung und Höflichkeit}

\begin{figure}[h]
  \centering
  \fbox{\rule{0pt}{3cm}\rule{8cm}{0pt}} % Platzhalter für KI-Bild
  \caption{Grüßen auf Deutsch}
\end{figure}

\section{Grüßen und Verabschieden}

\begin{lernkasten}[Wichtige Phrasen / Важливі фрази]
\begin{tabular}{p{5cm} p{5cm}}
\toprule
\textbf{Deutsch} & \textbf{Ukrainisch} \\
\midrule
Guten Morgen! & Доброго ранку! \\
Guten Tag! & Доброго дня! \\
Guten Abend! & Доброго вечора! \\
Auf Wiedersehen! & До побачення! \\
Tschüss! & Бувай! \\
Bitte & Будь ласка \\
Danke & Дякую \\
Danke schön! & Дуже дякую! \\
Entschuldigung & Вибачте \\
Wie geht es Ihnen? & Як ви почуваєтеся? \\
Gut, danke! & Добре, дякую! \\
\bottomrule
\end{tabular}
\end{lernkasten}

\section{Sich vorstellen}

\begin{lernkasten}[Ich stelle mich vor / Я представляюся]
\begin{tabular}{p{5cm} p{5cm}}
\toprule
\textbf{Deutsch} & \textbf{Ukrainisch} \\
\midrule
Wie heißen Sie? & Як вас звати? \\
Ich heiße \ldots & Мене звуть \ldots \\
Woher kommen Sie? & Звідки ви? \\
Ich komme aus der Ukraine. & Я з України. \\
Ich wohne in Mannheim. & Я живу у Мангаймі. \\
\bottomrule
\end{tabular}
\end{lernkasten}

\begin{uebungskasten}[Übung 1: Stell dich vor!]
Fülle die Lücken aus: \\[0.5em]
Ich heiße \underline{\hspace{4cm}}. \\
Ich komme aus \underline{\hspace{4cm}}. \\
Ich wohne in \underline{\hspace{4cm}}.
\end{uebungskasten}

% ============================================================
% KAPITEL 3: ALLTAG
% ============================================================
\chapter{Alltag — Supermarkt, Apotheke, Bank}

\section{Im Supermarkt}

\begin{lernkasten}[Einkaufen / Покупки]
\begin{tabular}{p{5cm} p{5cm}}
\toprule
\textbf{Deutsch} & \textbf{Ukrainisch} \\
\midrule
Wo ist der Supermarkt? & Де супермаркет? \\
Ich suche \ldots & Я шукаю \ldots \\
Was kostet das? & Скільки це коштує? \\
Das ist zu teuer. & Це занадто дорого. \\
Ich nehme das. & Я беру це. \\
Kann ich mit Karte zahlen? & Можна оплатити карткою? \\
Die Quittung, bitte! & Чек, будь ласка! \\
\bottomrule
\end{tabular}
\end{lernkasten}

\section{In der Apotheke}

\begin{lernkasten}[Apotheke / Аптека]
\begin{tabular}{p{5cm} p{5cm}}
\toprule
\textbf{Deutsch} & \textbf{Ukrainisch} \\
\midrule
Ich brauche ein Medikament. & Мені потрібні ліки. \\
Ich habe Kopfschmerzen. & У мене болить голова. \\
Ich habe Fieber. & У мене температура. \\
Haben Sie etwas gegen \ldots? & Чи є у вас щось від \ldots? \\
Ich brauche ein Rezept. & Мені потрібен рецепт. \\
\bottomrule
\end{tabular}
\end{lernkasten}

\section{Bei der Bank}

\begin{lernkasten}[Bank / Банк]
\begin{tabular}{p{5cm} p{5cm}}
\toprule
\textbf{Deutsch} & \textbf{Ukrainisch} \\
\midrule
Ich möchte ein Konto eröffnen. & Я хочу відкрити рахунок. \\
Ich brauche eine Überweisung. & Мені потрібен переказ. \\
Was sind die Gebühren? & Які комісії? \\
Ich habe meinen Ausweis. & У мене є посвідчення. \\
\bottomrule
\end{tabular}
\end{lernkasten}

% ============================================================
% KAPITEL 4: ARBEIT
% ============================================================
\chapter{Arbeit — Bewerbung und Vorstellungsgespräch}

\section{Bewerbungsunterlagen}

\begin{lernkasten}[Wichtige Dokumente / Важливі документи]
\begin{itemize}[nosep]
  \item \textbf{Lebenslauf (CV)} — Резюме
  \item \textbf{Anschreiben} — Супровідний лист
  \item \textbf{Zeugnisse} — Атестати / Дипломи
  \item \textbf{Referenzen} — Рекомендації
\end{itemize}
\end{lernkasten}

\section{Im Vorstellungsgespräch}

\begin{lernkasten}[Typische Fragen / Типові запитання]
\begin{tabular}{p{5cm} p{5cm}}
\toprule
\textbf{Deutsch} & \textbf{Ukrainisch} \\
\midrule
Erzählen Sie von sich. & Розкажіть про себе. \\
Was sind Ihre Stärken? & Які ваші сильні сторони? \\
Warum wollen Sie hier arbeiten? & Чому ви хочете тут працювати? \\
Wann können Sie anfangen? & Коли ви можете розпочати? \\
Was sind Ihre Gehaltsvorstellungen? & Які ваші зарплатні очікування? \\
\bottomrule
\end{tabular}
\end{lernkasten}

\begin{uebungskasten}[Übung: Bewerbungsgespräch]
Übe mit einem Partner. Einer fragt, der andere antwortet. \\
Wechselt nach 5 Minuten die Rollen!
\end{uebungskasten}

% ============================================================
% KAPITEL 5: ÜBUNGEN & QUIZ
% ============================================================
\chapter{Übungen \& Quiz}

\section{Quiz 1: Was sagt man?}

\begin{uebungskasten}[Wähle die richtige Antwort:]
\begin{enumerate}[nosep]
  \item Wie antwortest du auf ``Wie geht es Ihnen?''\\
    a) Danke schön! \quad b) Gut, danke! \quad c) Auf Wiedersehen!

  \item Was bedeutet ``Entschuldigung''?\\
    a) Дякую \quad b) Вибачте \quad c) Будь ласка

  \item Wie fragt man nach dem Preis?\\
    a) Was kostet das? \quad b) Ich nehme das. \quad c) Wo ist das?
\end{enumerate}
\end{uebungskasten}

\section{Quiz 2: Übersetze ins Deutsche}

\begin{uebungskasten}[Ukrainisch → Deutsch:]
\begin{enumerate}[nosep]
  \item Доброго ранку! → \underline{\hspace{5cm}}
  \item Дякую! → \underline{\hspace{5cm}}
  \item Де супермаркет? → \underline{\hspace{5cm}}
  \item У мене болить голова. → \underline{\hspace{5cm}}
\end{enumerate}
\end{uebungskasten}

\section{Antworten / Відповіді}

\begin{lernkasten}[Lösungen]
Quiz 1: 1b, 2b, 3a \\[0.3em]
Quiz 2:
\begin{enumerate}[nosep]
  \item Guten Morgen!
  \item Danke!
  \item Wo ist der Supermarkt?
  \item Ich habe Kopfschmerzen.
\end{enumerate}
\end{lernkasten}

% ============================================================
% KAPITEL 6: GLOSSAR
% ============================================================
\chapter{Glossar — Deutsch / Ukrainisch}

\begin{glossarkasten}[A–G]
\begin{multicols}{2}
\begin{description}[nosep,leftmargin=0pt]
  \item[Apotheke] Аптека
  \item[Arbeit] Робота
  \item[Arzt] Лікар
  \item[Auf Wiedersehen] До побачення
  \item[Ausweise] Документи
  \item[Bank] Банк
  \item[Bewerbung] Заява про прийом на роботу
  \item[Bitte] Будь ласка
  \item[Danke] Дякую
  \item[Deutsch] Німецька мова
  \item[Einkaufen] Покупки
  \item[Entschuldigung] Вибачте
  \item[Fieber] Температура
  \item[Geld] Гроші
  \item[Gespräch] Розмова
\end{description}
\end{multicols}
\end{glossarkasten}

\begin{glossarkasten}[H–Z]
\begin{multicols}{2}
\begin{description}[nosep,leftmargin=0pt]
  \item[Haus] Будинок
  \item[Hilfe] Допомога
  \item[Karte] Картка
  \item[Kopfschmerzen] Головний біль
  \item[Krankenhaus] Лікарня
  \item[Lebenslauf] Резюме
  \item[Mannheim] Мангайм
  \item[Medikament] Ліки
  \item[Rechnung] Рахунок
  \item[Rezept] Рецепт
  \item[Supermarkt] Супермаркет
  \item[Überweisung] Переказ
  \item[Wohnung] Квартира
  \item[Zeugnis] Диплом / Атестат
\end{description}
\end{multicols}
\end{glossarkasten}

% ============================================================
% SCHLUSS
% ============================================================
\chapter*{Viel Erfolg!}
\addcontentsline{toc}{chapter}{Viel Erfolg!}

\begin{center}
  {\color{gold}\rule{\linewidth}{2pt}}
  \vspace{1cm}

  {\Large\sffamily\bfseries Glückwunsch! Вітаємо!}

  \vspace{0.5cm}
  Du hast das Buch durchgearbeitet. \\
  Jetzt bist du bereit für den Alltag in Mannheim!

  \vspace{1cm}
  {\color{gold}\rule{\linewidth}{2pt}}

  \vspace{1cm}
  \textbf{Kontakt / Підтримка:}\\
  Telegram: @germancardsua\\

  \vspace{0.5cm}
  \textit{Deutsch lernen in Mannheim — Premium Edition}\\
  \textit{© 2026 German Cards Ukraine. Alle Rechte vorbehalten.}
\end{center}

\end{document}
